\documentclass[../../../../uem_mat_mei_q.tex]{subfiles}

\begin{document}
    集合$X$の要素の個数が有限であるとき，その個数を$n(X)$と表すことにする．%
    全体集合$U$および$U$の部分集合$A$，$B$がそれぞれ

    \myspaceb%
    $U=\{n\,|\,\text{$n$は15以下の自然数}\}$

    \myspaceb%
    $A=\{1,\,2,\,5,\,7,\,9,\,10,\,12,\,14,\,15\}$
    
    \myspaceb%
    $B=\{2,\,4,\,5,\,9,\,10,\,13,\,15\}$\\
    であるとき，

    \myegga{A}\myspacea$n(\overline{A}\cup\overline{B})$

    \myegga{B}\myspacea$n(A\cup B)$

    \myegga{C}\myspacea$n((A\cap B)\cup(\overline{A\cup B}))$
\end{document}
